\noindent\textbf{Elección del procedimiento.} La serie de esfuerzo normalizado tiene primer término $1-2e^{-\pi^2\nu t/h^2}$. Imponer 0,99 da $e^{-\pi^2\nu t/h^2}=0{,}005$, tomar logaritmo da $t=h^2\ln200/(\pi^2\nu)$. La magnitud del segundo término es $2e^{-4\pi^2\nu t/h^2}$ y permite comprobar la truncación.

\begin{enumerate}

\item Calcular la escala viscosa de la separación.

\[\fitmath{\mathrm{tnu}=\frac{{h}^{2}}{\nu}}\]

\[\fitmath{\mathrm{tnu}=\frac{{(0{,}01)}^{2}}{(1\times10^{-6})}\approx 100\,\mathrm{s}}\]

\item Evaluar el término omitido a 50 s.

\[\fitmath{\mathrm{term}_{2}=2\cdot \exp\left(\frac{-4\cdot {\pi}^{2}\cdot \nu\cdot \mathrm{tcheck}}{{h}^{2}}\right)}\]

\[\fitmath{\mathrm{term}_{2}=2\cdot \exp\left(\frac{-4\cdot {\pi}^{2}\cdot (1\times10^{-6})\cdot (50)}{{(0{,}01)}^{2}}\right)\approx 5{,}350576\times10^{-9}\,}\]

\item Despejar el instante del criterio del 99 por ciento en el primer modo.

\[\fitmath{t_{99}=\frac{{h}^{2}\cdot \ln\left(200\right)}{{\pi}^{2}\cdot \nu}}\]

\[\fitmath{t_{99}=\frac{{(0{,}01)}^{2}\cdot \ln\left(200\right)}{{\pi}^{2}\cdot (1\times10^{-6})}\approx 53{,}68318\,\mathrm{s}}\]

\end{enumerate}

\noindent\textbf{Contraste y comparación.} El primer modo da 53,68 s y el segundo término es diminuto. Un valor numérico de 55,1 s no se justifica por esa truncación: depende de discretización, muestreo temporal o criterio de detección, que deben refinarse antes de atribuirle mayor exactitud.
